\section{一次长任务给我的实际认识}
我在这次实验里最明确的修改，是不再要求GPT先写出所有后续实验的详细顺序。我仍要保留教师范围、数据语义、指标、预算和最终确认这些边界，但允许Codex在边界内根据真实结果安排必要的单项诊断、有限组合和回退。这种模式是否有效，不能靠“单Goal”三个字证明，而要看真实实验、错误修复和独立检查的记录。

实际执行留下了相互支持也相互限制的结论。TBG阈值提高可以降低平均延迟，却未保持最大延迟的优势；KM混合初始化在已测输入上改善前沿，但扩大预算的附加收益很小；Q-KM在不同车队规模下有正有负；主动调度的gross增益要与真实产生的空驶量和成本假设一起读。这样的结果使我更明确：\textbf{完成算法、支持效果、适合采用和整体交付，是四种不同判断。}

我也看到了工作流自身需要保持的边界。Codex内部审查是真实工程证据，却不能替代GPT读取固定GitHub版本的二审；Agent自动提出和执行的候选，不能在Process里写成我逐项追问；录屏提供的UI是人的真实要求和回应，不能把后来形成的专业术语倒写成最初的自然语言。正因为这些边界被记录下来，报告才能同时展示我的判断与自动研究的实际贡献。

下一轮项目工作可以继承这种做法，但不能直接继承本次的批次长度、车队数量、参数扫描或长任务授权。具体研究合同仍应由我在任务开始时确认；在工作流本身没有必要的新增价值时，也不应为显示复杂性而制造新的Agent、图表和流程层。